%% Example CV - Henrique Cordeiro da Silva
%% Comprehensive overview of all competencies, experience, and achievements.
%% Use as a master reference when tailoring role-specific CVs.
%%
%% Compile with: lualatex main_example.tex
%% (pdflatex often fails on modern MiKTeX with fontawesome5 font-expansion errors.)

\documentclass[11pt,a4paper,sans]{moderncv}
\moderncvstyle{banking}
\moderncvcolor{blue}

% Force both first and last name AND section headings to render in moderncv
% blue (color1). Default banking on lualatex+MiKTeX leaves these black, which
% looks inconsistent with the rest of the blue accent scheme.
\renewcommand*{\firstnamestyle}[1]{{\fontsize{34}{36}\bfseries\upshape\color{color1}#1}}
\renewcommand*{\lastnamestyle}[1]{{\fontsize{34}{36}\bfseries\upshape\color{color1}#1}}
\renewcommand*{\sectionstyle}[1]{{\sectionfont\color{color1}#1}}

\usepackage[utf8]{inputenc}
\usepackage{hyperref}
\hypersetup{
    colorlinks=true,
    linkcolor=blue,
    filecolor=magenta,
    urlcolor=blue,
    pdftitle={Henrique Cordeiro da Silva - CV},
    pdfpagemode=FullScreen,
}
\usepackage[scale=0.80]{geometry}
\usepackage{import}
\usepackage{needspace}

% personal data
\name{Henrique}{Cordeiro da Silva}
\address{Hortolândia, SP, Brasil}{}{}
\phone[mobile]{+55 19 99545-4391}
\email{henriquecordeiro054@gmail.com}
\extrainfo{\href{https://linkedin.com/in/henriqucordeiro}{LinkedIn}, \href{https://github.com/HenriqueCDS}{GitHub}, \href{https://portifolio-henriqueocordeiro.vercel.app}{Portfólio}}

\begin{document}

\makecvtitle

% ============================================================
%     PROFILE STATEMENT
% ============================================================

\vspace{6pt}
\small{Desenvolvedor Backend com atuação em Python e Java, especializado em APIs REST, integração de sistemas e pipelines de ETL. Sustento o Canvas LMS na PUC-Campinas (16.500 usuários, 150.000+ registros acadêmicos por ciclo letivo) e sou pós-graduado em Ciência de Dados e Machine Learning, com experiência prática construindo agentes de IA (RAG) em produção.}

% ============================================================
%     CORE COMPETENCIES
% ============================================================

\section{Competências Principais}
\vspace{1pt}
\begin{itemize}

\item \textbf{Backend e APIs}: Python, Java, APIs REST, Spring Boot, FastAPI, .NET, Node.js, microsserviços, integração de sistemas, mensageria (RabbitMQ)

\item \textbf{Dados e ETL}: Pandas, NumPy, modelagem de dados, SQL avançado, MySQL, PostgreSQL, SQL Server, MongoDB, pgvector

\item \textbf{IA e Machine Learning}: LangChain, RAG (Retrieval-Augmented Generation), embeddings, API Gemini, Scikit-learn, redes neurais

\item \textbf{DevOps e Ferramentas}: Git, GitHub, Docker, Linux, CI/CD, Maven, Postman

\item \textbf{Frontend}: React, JavaScript, TypeScript, HTML, CSS

\end{itemize}

% ============================================================
%     PROFESSIONAL EXPERIENCE
% ============================================================

\section{Experiência Profissional}
\vspace{3pt}
\begin{itemize}

% --- Most Recent Role ---
\needspace{5\baselineskip}
\item{\cventry{jul/2024--atual}{Analista de Suporte / Desenvolvedor}{PUC-Campinas}{Campinas, SP}{}{\vspace{1pt}
\begin{itemize}
    \item Garanto a disponibilidade e a estabilidade do Canvas LMS para 16.500 alunos e docentes, atuando na análise de logs, diagnóstico de incidentes em produção e monitoramento de integrações
    \item Desenvolvo automações em Python integradas à API REST do Canvas para criação de salas, sincronização de matrículas e auditoria de acessos
    \item Asseguro a integridade de mais de 150.000 registros acadêmicos por ciclo letivo, implementando pipelines de ETL em Python e Pandas
    \item Elimino o retrabalho de conferência manual de matrículas com rotinas automatizadas em Jupyter Notebook
    \item Mantenho e evoluo uma aplicação .NET em produção, corrigindo falhas e implementando melhorias solicitadas pelo negócio
\end{itemize}}}

\vspace{3pt}

% --- Previous Role ---
\needspace{5\baselineskip}
\item{\cventry{mar/2022--mar/2024}{Desenvolvedor Full Stack (Estágio)}{FUNCAMP}{Campinas, SP}{}{\vspace{1pt}
\begin{itemize}
    \item Automatizei processos de negócio entre sistemas internos e o ERP NetSuite com scripts de integração em Java
    \item Construí uma API REST para armazenamento e distribuição centralizada do banco de questões da plataforma Edukas
    \item Entreguei telas e funcionalidades em PHP e JavaScript, atuando em time ágil de 3 a 6 pessoas
\end{itemize}}}

\end{itemize}

% ============================================================
%     EDUCATION
% ============================================================

\section{Formação Acadêmica}
\vspace{1pt}
\begin{itemize}

\needspace{4\baselineskip}
\item{\cventry{jun/2025--jun/2026}{Pós-graduação em Ciência de Dados e Machine Learning}{PUC-Campinas}{Campinas, SP}{}{\vspace{1pt}
Concluído.
}}

\vspace{3pt}

\needspace{4\baselineskip}
\item{\cventry{fev/2022--jun/2024}{Tecnólogo em Análise e Desenvolvimento de Sistemas}{Centro Universitário FAVIP Wyden}{Campinas, SP}{}{\vspace{1pt}
Concluído.
}}

\vspace{3pt}

\needspace{4\baselineskip}
\item{\cventry{fev/2018--dez/2021}{Técnico em Informática (Ensino Médio Integrado)}{IFSP Campus Campinas}{Campinas, SP}{}{\vspace{1pt}
Concluído.
}}

\end{itemize}

% ============================================================
%     LANGUAGES
% ============================================================

\section{Idiomas}
\vspace{1pt}
\begin{itemize}
\item Português (nativo), Inglês (básico -- leitura de documentação técnica).
\end{itemize}

% ============================================================
%     PROJECTS
% ============================================================

\section{Projetos}
\vspace{1pt}
\begin{itemize}
\item \textbf{ia-agent-puc-digital} -- agente de RAG em Python (FastAPI, LangChain) para suporte acadêmico, com embeddings locais, busca por similaridade em pgvector, integração com a API Gemini e 25 testes automatizados (pytest); containerizado com Docker. \href{https://github.com/HenriqueCDS}{github.com/HenriqueCDS}
\item \textbf{proposta-app-mensageria} -- sistema de análise de crédito em microsserviços com Java, Spring Boot, RabbitMQ, PostgreSQL, AWS SNS e Docker.
\item \textbf{Rest\_Api\_Questoes} -- API REST em Node.js, Express e Sequelize, com testes em Jest e persistência em MySQL.
\end{itemize}

% ============================================================
%     REFERENCES
% ============================================================

\section{Referências}
\vspace{1pt}
\begin{itemize}
\item{Disponíveis mediante solicitação.}
\end{itemize}

\end{document}
